\par\Needspace{6\baselineskip}\medskip
\hypertarget{native-084-article065}{}\pdfbookmark[2]{65.}{article-065}
\begin{center}65.\end{center}
Etiamsi in eo casu, ubi ex elementis datis paucissimi tantum loci geocentrici sunt determinandi, omnia artificia supra tradita, per quae ab anomalia excentrica statim ad longitudinem et latitudinem geocentricam, vel adeo ad rectascensionem et declinationem, transire licet, in usum vocare vix operae pretium sit, quoniam compendia inde demanantia a multitudine quantitatum auxiliarium antea computandarum absorberentur: semper tamen contractio reductionis ad eclipticam cum calculo longitudinis et latitudinis geocentricae lucrum haud spernendum praestabit. Si enim pro plano coordinatarum $z$ assumitur ipsa ecliptica, poli autem planorum coordinatarum $x$, $y$ collocantur in longitudine $\AscendingNode$, $90^0+\AscendingNode$, coordinatae facillime absque ulla quantitatum auxiliarium necessitate determinantur. Habetur scilicet
\SourceLeaf{085}
\[
\begin{array}{@{}l|l|l@{}}
x=r\cos u & X=R'\cos(L-\AscendingNode) & x-X=\Delta'\cos(l-\AscendingNode)\\[.4em]
y=r\cos i\sin u & Y=R'\sin(L-\AscendingNode) & y-Y=\Delta'\sin(l-\AscendingNode)\\[.4em]
z=r\sin i\sin u & Z=R'\operatorname{tang}B & z-Z=\Delta'\operatorname{tang}b.
\end{array}
\]
\noindent Quoties $B=0$, est $R'=R$, $Z=0$. Secundum has formulas exemplum nostrum numeris sequentibus absolvitur: $L-\AscendingNode=213^0\,12^{\prime}\,0^{\prime\prime}{,}32$
\begin{center}
\begin{tabular}{@{}ll@{}}
\makebox[13em][l]{$\log r$\hfill .\,.\,.\,.\,.\,.\,.\,.\,.\,.\,.} & $0{,}325\,9877$\\[.3em]
\makebox[13em][l]{$\log\cos u$\hfill .\,.\,.\,.\,.\,.\,.\,.} & $9{,}982\,4141_n$\\[.3em]
\makebox[13em][l]{$\log\sin u$\hfill .\,.\,.\,.\,.\,.\,.\,.} & $9{,}445\,4714_n$\\[.3em]
\hline
\makebox[13em][l]{$\log x$\hfill .\,.\,.\,.\,.\,.\,.\,.\,.\,.\,.} & $0{,}308\,4018_n$\\[.3em]
\hline
\makebox[13em][l]{$\log r\sin u$\hfill .\,.\,.\,.\,.\,.\,.} & $9{,}771\,4591_n$\\[.3em]
\makebox[13em][l]{$\log\cos i$\hfill .\,.\,.\,.\,.\,.\,.\,.} & $9{,}988\,5266$\\[.3em]
\makebox[13em][l]{$\log\sin i$\hfill .\,.\,.\,.\,.\,.\,.\,.} & $9{,}355\,7570$\\[.3em]
\hline
\makebox[13em][l]{$\log y$\hfill .\,.\,.\,.\,.\,.\,.\,.\,.\,.\,.} & $9{,}759\,9857_n$\\[.3em]
\makebox[13em][l]{$\log z$\hfill .\,.\,.\,.\,.\,.\,.\,.\,.\,.\,.} & $9{,}127\,2161_n$
\end{tabular}
\end{center}
\begin{center}
\begin{tabular}{@{}ll@{}}
\makebox[13em][l]{$\log R'$\hfill .\,.\,.\,.\,.\,.\,.\,.\,.\,.} & $9{,}998\,0979$\\[.3em]
\makebox[13em][l]{$\log\cos(L-\AscendingNode)$\hfill .\,.\,.} & $9{,}922\,6027_n$\\[.3em]
\makebox[13em][l]{$\log\sin(L-\AscendingNode)$\hfill .\,.\,.} & $9{,}738\,4353_n$\\[.3em]
\hline
\makebox[13em][l]{$\log X$\hfill .\,.\,.\,.\,.\,.\,.\,.\,.\,.} & $9{,}920\,7006_n$\\[.3em]
\makebox[13em][l]{$\log Y$\hfill .\,.\,.\,.\,.\,.\,.\,.\,.\,.} & $9{,}736\,5332_n$\\[.3em]
$Z=$ & $0$
\end{tabular}
\end{center}
\noindent Hinc fit
\begin{center}
\begin{tabular}{@{}ll@{}}
\makebox[13em][l]{$\log(x-X)$\hfill .\,.\,.\,.\,.\,.} & $0{,}079\,5906_n$\\[.3em]
\makebox[13em][l]{$\log(y-Y)$\hfill .\,.\,.\,.\,.\,.} & $8{,}480\,7165_n$\\[.3em]
unde $(l-\AscendingNode)=$ & $181^0\,26^{\prime}\,33^{\prime\prime}{,}49$\\[.3em]
\makebox[13em][l]{$\log\Delta'$\hfill .\,.\,.\,.\,.\,.\,.\,.\,.\,.} & $0{,}079\,7283$\\[.3em]
\makebox[13em][l]{$\log\operatorname{tang}b$\hfill .\,.\,.\,.\,.\,.\,.} & $9{,}047\,4878_n$
\end{tabular}
\end{center}
\begin{center}
\begin{tabular}{@{}rl@{}}
$l=$ & $352^0\,34^{\prime}\,22^{\prime\prime}{,}22$\\[.4em]
$b=$ & $-6\quad21\quad55{,}06$
\end{tabular}
\end{center}
